%
%      Brno University of Technology
%      Faculty of Information Technology
%
%      MSc Thesis	2009/2010
%
%      Design of S-Boxes Using Genetic Algorithms
%
%
%      Bedrich Hovorka
%
Cryptography \cite{yokohama, bethany} deals primarily
with concealing information and verifying its origin,
i.e., transforming data into a form that
is unusable for unwanted parties, which will hereafter
 in the text be collectively referred to as the \uv{attacker} for simplicity, even though they may not attack directly
 and there may be multiple of them. Authorized parties have, in contrast to unauthorized ones, special knowledge.
 In the past, this knowledge included the encryption procedure itself, but nowadays this is abandoned and
 the knowledge is only the key. It is assumed that the attacker also knows the algorithm. Furthermore, to non-concealed
 information, additional information about its origin can be added,
 so-called digital signature or timestamp, in such a way that
 it is practically impossible to verify with this signature any information other than the one that identically corresponds to the one at signing.
Such a transformation should solve the main security problems:
\begin{description}
  \item[confidentiality] the attacker must not obtain any useful information from the encrypted
  communication channel
  \item[authentication] verification of the sender's identity
  \item[integrity] the attacker cannot forge data or parts thereof and present them to the recipient as if they originated
    from a specific sender -- i.e., also change or modify data originating from the legitimate sender
  \item[non-repudiation] the sender cannot deny that a message created and signed by them was sent by them
\end{description}

This transformation can be simplified into a mathematical expression $C = F(K_1,P)$, where $P$
is plaintext, $C$ is ciphertext,
and $K$ is the key, which represents the possibility to set how the cipher transforms the text.
For the reverse process called decryption $P = F^{-1}(K_2,C)$, the same key need not be used.
The process is also illustrated in Figure \ref{cryptoIntro}.
The longer the key, the more difficult it is to perform a brute force attack to obtain it.
The function $F$ is hereafter referred to as a cipher or also a cryptosystem.

\insertimage[0.8\textwidth]{cryptoIntro}{Illustration of the information concealment process}

Modern cryptography is divided into two main categories (according to the number of keys):
\begin{description}
  \item[symmetric] A symmetric cryptosystem is one
    where the decryption key can be easily derived from the encryption key.
    Usually there exists only one key, the so-called \uv{secret key},
    which is used for both encryption and decryption. The problem with using
    this system is key distribution.
  \item[asymmetric] An asymmetric system is one where the decryption
  key cannot be practically derived from the encryption key.
  Usually there exists a pair of keys, which is generated by the recipient, one of the keys in the pair is
  \uv{public} and in the information concealment process it is used by the sender for encryption.
  The second key is \uv{private} (no one except the recipient may know it),
  with which the recipient decrypts the received message. The public key can be, thanks to the fact that
  the private key cannot be easily derived from it, exposed anywhere.
  Key distribution is solved at the cost of
  computational complexity of encryption.
\end{description}
Both approaches are usually combined in practice such that
the sender first generates a random secret key $K_{rand}$,
which is encrypted with an asymmetric cipher using the recipient's public key
$$C_1 = F_{asym}(K_{public_{recv}}, K_{rand})$$
The actual message $P$ is then encrypted symmetrically with the key $K_{rand}$:
$C_2 = F_{sym}(K_{rand}, P)$, because symmetric encryption is considerably faster.
Finally, the sender sends data in the form $C_1, C_2$ to the recipient (expressed in simplified form).
By block is hereafter meant any sequence of consecutive bits.
By encryption block, a \uv{box} that transforms a block of plaintext into a block of ciphertext.
Individual blocks will be further represented as signals that
indicate connections between individual parts of the algorithm and information flow.

In the previous paragraphs, the procedure for concealing information was illustrated.
For digital signatures, an asymmetric cipher is used with the sender's keys reversed.
When signing, the sender uses their private key and the recipient decrypts it
using the sender's public key, thereby verifying the origin of the message.
The sender must have two different key pairs, one pair for encryption
and another pair for signing, so that an attacker cannot substitute another piece of information for signing
that is intended for the sender and which they would consequently decrypt.
Moreover, because asymmetric encryption is a computationally expensive operation,
when signing a message, a fingerprint (so-called \uv{hash}) is first computed and it
is then asymmetrically encrypted with the sender's private key and attached to the message.
The recipient verifies the message by first computing the \uv{hash} of the received message
and simultaneously decrypting the received \uv{hash} using the sender's public key,
then comparing both \uv{hashes}.
Hash algorithms are another group of cryptographic algorithms.
A hash algorithm thus performs a transformation of an arbitrarily long message into a fixed
length, the so-called \uv{hash}.
This \uv{hash} should be unique to the given message. Two different messages, even if they differ
by one bit, must not have the same \uv{hash}.

\begin{table}[t]
\begin{center}
  \begin{tabular}{| l | c | c|}
  \hline
  Type & Abbreviation & key bit count \\ \hline
  symmetric & AES & 128, 192, 256\\ \hline
  ~ & DES & 56, 168 \\ \hline
  ~ & RC4 & 40 to 256 \\ \hline
  asymmetric & RSA & \dots, 1024, 2048, 4096, \dots \\ \hline
  ~ & DSA & 1024, 2048, 3072\\ \hline
  hash & & \\ \hline
  ~ & MD5 & 128\\ \hline
  ~ & SHA1 & 160\\ \hline
  ~ & SHA2 & 224, 256, 384, 512\\ \hline
  \end{tabular}
  \caption{Examples of the most well-known cryptographic algorithms}
  \label{cipherExamples}
\end{center}
\end{table}
Table \ref{cipherExamples} lists several of the most commonly used cryptographic algorithms and their
most commonly used key lengths. For asymmetric algorithms, the key length increases over the years
much more than for symmetric and hash algorithms. Developing a new symmetric cipher is much simpler than
developing an asymmetric one. For asymmetric ciphers, it is necessary to first find a mathematical principle, therefore their security
is more reflected in the key size. Symmetric ciphers with keys smaller than 80 bits are considered weak.
For asymmetric ciphers, this threshold is already at 1024 bits. Hash algorithms MD5 and SHA1 are also considered
weak.
The 168-bit key for DES is called TripleDES,
this key is divided into three 56-bit keys and during encryption goes through three phases of DES
$C = F(K_3,(F^{-1}(K_2(F(K_1, P)))$, the middle one is decryption. Only keyless hash algorithms are listed.
NIST\footnote{National Institute of Standards and Technology USA} announced a competition for a new SHA3 standard,
which is currently in the 2nd round, among the advancing candidates are for example
CubeHash and Luffa (\cite{specLuffa}).

\section{Symmetric cryptography}
Symmetric ciphers are divided into block and stream ciphers.
Stream ciphers take messages byte by byte, or bit by bit or character by character. Block ciphers
take messages in larger blocks (64 to 256 bits). Larger block size prevents statistical
and dictionary attacks. For an $n$-bit block there are $2^n$ possible input and output combinations and
there exist $2^{n}!$ possible mappings.
An example of a block cipher is DES, which encrypts in 64-bit blocks. An example of a stream cipher
is RC4. A block cipher can become a stream cipher if it is connected in a certain mode (see below).

\subsubsection{Block cipher modes}
Divided message blocks can be processed by any symmetric block cipher
in the following ways. If the attacker has available corresponding pairs of plaintext and ciphertext blocks,
some modes may be weaker. In some modes, an initialization vector is added to the key, which is
used for initialization of the first block encryption, as additional knowledge necessary for decryption.
\begin{description}
  \item[Electronic Codebook (ECB)] The simplest, each $i$-th block of the message is directly the input
  of the encryption block and the output
  is directly the $i$-th block of ciphertext. Cryptanalysis of this mode is simpler,
  because the attacker can find the same block in the ciphertext again, for which
  they already have the corresponding plaintext in a dictionary that they continuously create.
  This also violates data integrity, as the attacker can substitute a different block of ciphertext.
  \item[Cipher Block Chaining (CBC)] Blocks are chained. Unlike ECB, each block is XORed
  with the previous output block before being used as input to the encryption block.
  The first block is XORed with the initialization vector. Illustrated in Figure \ref{cbc}.
  \item[Cipher Feedback (CFB)] Unlike CBC mode, the plaintext block is XORed with the output
  of the cipher block and the result is also fed to the input of the following encryption block. The input of the first block is
  the initialization vector.
  \item[Output Feedback (OFB)] Unlike CFB mode, the output of the previous block is sent to the input of the following block
  before XORing with the plaintext.
\end{description}

\insertimage[0.8\textwidth]{cbc}{CBC mode \cite{wikiCBCMAC}}

In addition to these, there is also the so-called \uv{counter mode}, where the block cipher is used to generate a pseudorandom
sequence, such that counter values are fed to the input of the block cipher. The advantage is that each block can
be encrypted independently of the others. The plaintext is XORed with the output of the block cipher.
The counter value must obviously not repeat for the same key.
OFB mode can also be referred to as a stream cipher and CFB mode as a self-synchronizing stream cipher.

\pagebreak % ???

\section{Internal structure of the encryption block}
The encryption of one block is divided into various phases according to the type of algorithm.
Many modern ciphers are based on the Feistel substitution/permutation network,
i.e., they are a combination of these transformations, occurring in several consecutive rounds:
\begin{description}
  \item[substitution] A certain block of code is replaced by another from a set of possible combinations.
  \item[transposition/permutation] Works with several blocks whose order is changed.
  \item[mixing with subkey] Subkeys are generated from the key for individual rounds,
  which are then \uv{mixed in}, most often with XOR.
\end{description}
\insertimage[0.7\textwidth]{spn}{Simple substitution/permutation network illustrating an encryption algorithm \cite{tutorialLADA}}
Figure \ref{spn} illustrates a simple multi-round
cipher.
Substitution is usually abstracted using so-called \uv{substitution boxes}, abbreviated S-boxes
(elements $S_{kn}$ from Figure \ref{spn}). S-boxes are represented as
\uv{boxes} $S_{kn}$, into which signals enter
and from which signals exit.
Permutation is represented using signals.
In each round
and at each position, S-boxes with different mappings are used.
In the illustrated figure, these are S-boxes of type $4 \times 4$.
In this work, we will be interested in their internals and outputs.
The principles of S-boxes are best described in symmetric ciphers, although their use
in other types of algorithms is not excluded.
Their issues can be best shown precisely in symmetric ciphers, because
they themselves have the character of a symmetric block cipher, albeit in a smaller form.

This simple network by itself, however, does not guarantee the possibility of both encryption and decryption using
the same structure. For the process to be reversible, that is, for it to be possible to encrypt/decrypt
in multiple rounds using the same structure,
Feistel \cite{wikiFeistel} proposed the structure in Figure \ref{Feistel}.

\insertimage[0.35\textwidth]{Feistel}{Structure of one round of a Feistel cipher \cite{wikiFeistel}}

The input text is divided into two halves. In each round, the right half is the input to both the encryption
function $F$ and the left half of the round output. The left half of the input is XORed with the output of the encryption
function and the result is the right half of the output. The encryption function is configured by the corresponding
subkey. For decryption, the subkeys are used in reverse order.
The responsibility for security is transferred to function $F$, while reversibility
is handled by the Feistel structure.
When for encryption in each round the following holds (calculation of the next round):
\begin{eqnarray}
L_{i+1} & = & R_i \\
R_{i+1} & = & L_i \oplus F(K_i, R_i)
\end{eqnarray}
then for decryption in each round the following holds:
\begin{eqnarray}
R_i & = & L_{i+1} \\
L_i & = & R_{i+1} \oplus F(K_{i+1}, R_{i+1})
\end{eqnarray}
$L_i$ resp. $R_i$ resp. $K_i$ is the left input resp. right input resp. subkey of the round.
In function $F$, substitution, permutation and transposition can be performed. It can contain several S-boxes.
This structure was adopted by many ciphers, e.g., DES or Twofish. They differ primarily in function $F$,
subkey generation, and the size of the processed block. Figure \ref{DESboxes} shows
the function in the DES cipher. $E$ is an expansion from 32 bits to 48 bits. This is followed by mixing with a 48-bit subkey.
The result is divided among 8 S-boxes of type $6 \times 4$. Their main task is
the non-linear mapping of their input to output (explained later). The outputs of the S-boxes are merged into 32 bits and
permuted by function $P$.

\insertimage[0.7\textwidth]{DESboxes}{Structure of the Feistel function in DES \cite{yokohama}}
In subsection \ref{secSboxUsing} (after explaining the following concepts), some S-boxes
are described in more detail.

\section{Achieving security}
When designing a cryptosystem, these questions should be asked:
\begin{itemize}
  \item How can a cryptologist prove the security of a cryptosystem?
  \item What statistical properties does the structure of the cryptosystem have?
  \item How can these properties be determined efficiently?
\end{itemize}
The first question cannot be answered fully satisfactorily, because with increasing computer performance
and the efforts of cryptanalysts, the discovery of a cipher weakness and exploitation of this discovery
in an attack is not excluded. Therefore, a cipher can be designed with the goal of resistance to all known attacks, while
it is not practically possible to completely validate the entire cipher. The strength of cipher components can be measured using
various criteria.

The answer to the second question is that \cite{yokohama} a cipher must not have a small number of
possible pairs of plaintexts and corresponding ciphertexts for different lengths from the entire text.
Any set of output bits should not depend on any proper subset
of the set of all input bits for a fixed key.
Any set of output bits should not depend on any proper subset
of the set of all key bits for a fixed input text. Among usable statistical tests,
one can mention for example the $\chi^2$-test or the Kolmogorov-Smirnov test. However, it is not practically
possible to go through all possible keys, let alone all possible texts. Therefore, several
random keys and random texts are generated. The output should have characteristics
of a random sequence. For verification, for example, \uv{frequency analysis} is used,
which verifies the occurrences of values at different positions and their dependencies.
It is clear that an ideally random output cannot be practically created.

Individual cipher components must first meet security properties. Different components are
responsible for different properties of the entire cipher.
When constructing individual components, thanks to the small number of inputs and outputs,
it can be undemanding to verify all these properties for all possible combinations.
Ways to verify the resistance of a component also include performing a known attack that
exploits weaknesses caused by poor characteristics of the cipher component.
The two most important attacks on encryption algorithms, which are used among other things to verify
the randomness of S-box outputs, were created to attack the DES algorithm.

\section{Linear cryptanalysis}
Linear cryptanalysis \cite{tutorialLADA} attempts to exploit biased probability of occurrence of linear expressions
involving bits of plaintext, ciphertext, and (sub)keys (this in the general case,
not considered further in this work).
This is a so-called \uv{known plaintext attack}: this means that it is assumed that
the attacker has available a set of plaintexts and their corresponding ciphertexts.
Nevertheless, the attacker does not have the option to choose which plaintext (and
corresponding ciphertext) is available. In many applications and scenarios, it is considered sufficient that
this set of texts is random.
In other words, linear cryptanalysis attempts from a random set of plaintexts and corresponding ciphertexts to
approximate the encryption algorithm with various
linear expressions and use the most (un)likely of them to gradually obtain the key.


Linearity is related to the mod-2 bit operation (e.g.,
exclusive sum denoted by \uv{$\oplus$}). And a linear expression is in the form:
\begin{equation}
  a_1 \cdot X_{i_1} \oplus a_2 \cdot X_{i_2} \oplus \dots \oplus a_u \cdot X_{i_u} \oplus
  b_1 \cdot Y_{j_1} \oplus a_2 \cdot Y_{j_2} \oplus \dots \oplus a_v \cdot Y_{j_v} = 0
  \label{eqLinearAprox}
\end{equation}

where $X_i$ represents the $i$-th input plaintext $X_i = [X_{i_1}, {X_{i_2}}, \dots]$
and $Y_j$ represents the $j$-th output ciphertext
$Y_i = [Y_{j_1}, Y_{j_2}, \dots]$.\footnote{Both input and output text are defined as bit vectors.}
This equation represents the \uv{sum} of exclusive sums of input and output bits that are
selected\footnote{\uv{Turned on} using bit conjunction denoted as \uv{$\cdot$}}
by bit vectors $a$ (for input) and $b$ (for output).


The concept of linear cryptanalysis is to determine expressions in the form of equation \ref{eqLinearAprox} that have
high or conversely low probability. If there is a tendency toward such probabilities,
this evidently means a weak ability of the cipher to produce output with characteristics
of randomly generated output.
Let us have randomly generated bits $u + v$ and if we use them in equation \ref{eqLinearAprox},
the probability $p_L$ should be exactly $\frac{1}{2}$. The deviation from probability
$p_L = \frac{1}{2}$ we call \uv{bias} denoted $\epsilon_i$, where $i$ relates to the expression.
The higher the magnitude of this deviation, $|\epsilon_i = p_L - \frac{1}{2}|$, the easier it is to break the cipher.

Note that $p_L = 1$ means that the linear expression \ref{eqLinearAprox} is an exact representation
of the cipher and consequently the cipher is catastrophically weak. If $p_L = 0$, then the expression represents
an affine relationship, also a sign of great weakness. For mod-2 arithmetic,
an affine function is simply the complement of a linear one. Both
linear resp. affine approximations, indicated by probabilities $p_L > \frac{1}{2}$ resp.  $p_L < \frac{1}{2}$,
have equal weight in linear cryptanalysis and hereafter in this text the term \uv{linear} will mean
both relationships.

\pagebreak
\section{Differential cryptanalysis}
Differential cryptanalysis \cite{tutorialLADA} exploits
the high probability of joint occurrence of differences between two ciphertexts
at the output with differences between two plaintexts at the input. For example, let us have
a system with input $X = [X_1 X_2 \dots X_n]$ and output $Y = [Y_1 Y_2 \dots Y_n]$.
Let two inputs of the system be $X'$ and $X''$ with corresponding outputs $Y'$ and $Y''$. The input
difference is given by $\Delta X = X' \oplus X''$ where \uv{$\oplus$} represents
the bitwise exclusive sum of $n$-bit vectors, and hence
\begin{equation}
  \Delta X = [\Delta X_1 \Delta X_2 \dots \Delta X_n ]
\end{equation}

where $\Delta X_i = X'_i \oplus X''_i$ where $X'_i$ resp. $X''_i$ represents the $i$-th bit of $X'$ resp. $X''$.
Similarly for outputs.

In an ideal cipher\footnote{i.e., one that has random properties at the output}
the probability that a partial output difference $\Delta Y$ occurs
for a given partial input difference $\Delta X$ is $p_D = \frac{1}{2}^n$, where $n$ is the number of bits in $X$.
Differential cryptanalysis is applicable in cases where $\Delta Y$
occurs together with $\Delta X$ with very high $p_D$ (e.g., much greater than
$\frac{1}{2}^n$). The pair $(\Delta X, \Delta Y)$ is called a \emph{differential}.

Differential cryptanalysis is a so-called \uv{chosen plaintext attack},
which means that the attacker can choose the input and examine the output in an attempt to derive the key.
In differential cryptanalysis, the attacker can choose a pair of inputs $X'$ and $X''$
belonging to $\Delta X$ with the knowledge that $\Delta X$ occurs with the highest probability
for $\Delta Y$.
